\section{Introduction}

Energy Price Forecasting is a highly sought after technique that adds value to companies that benefit from load balancing and management. With the further decarbonisation of the energy grid, renewable energy introduces
a strong temporal characteristic to the challenge; one that Neural Networks, particularly advancements in RNN\cite{Aksu2026,Xu2025,Shao2021,Fang2023,Manero2018}, Transformer\cite{Kong2025,Wang2026,Oliveira2023,9756020} and ensemble\cite{math11030676,CHEN2025126302,Hussan2026-dy} technology can leverage to improve the accuracy of predictions. As feature extraction mechanisms grow more robust\cite{CHEN2025126302,Aksu2026,Xu2025,4494757,Shao2021},
and models involve complex sets of features like weather\cite{9788043, Wang2026}, location, and energy consumption\cite{4494757}, the demand for better, more accurace, and more farreaching models is growing withing the industry.

One of the industry players with a keen interest in the technology is OpenRemote. The company provides a platform where entities with multiple energy-demanding components can manage their resources. One of the microservices
that OpenRemote provides is an Energy Price Forecasting (EPF) model hosted by Prophet, that does not employ Neural Networks for its predictions\cite{Taylor2017}. Their current service is only able to offer 24h
predictions ahead of time, which causes constraint for schedueled load balancing, and limits the amount of saving a client could make. OpenRemote argues that a 48h prediction window would be more benefitial. Furthemore, the Prophet model is more accurate around the peaks, but 
relatively innacurate anywhere else. This is due to the limitations of the modeling techniques employed which, although easy to configure and faster than a Neural Network to train and generate forecastings, struggles 
with the complexities of multi-feature temporal and spacial data that RNNs and Transformers capture with greater ease.

This paper will present, out of a collection of methods collected via a thourough literature review, the results of experiments conducted with the intention of comparing
different combinations of methods - both preprocesing, and neural network architectures. The end result if the analysis and thourough review of the results. It will also provide
advice on how to implement the techniqeus reviewed throughut this research.